\section{Translation Schemes}
\label{sec:ch14-scheme}

\subsection{Definisi}

Translation scheme adalah tata bahasa dengan potongan program
semantik yang disisipkan di tengah ruas kanan. Berbeda dari SDD yang
mendeskripsikan aturan, translation scheme mendeskripsikan urutan
eksekusi.

\begin{definisi}[Translation scheme]
Translation scheme adalah tata bahasa dengan aksi semantik yang
ditempatkan pada posisi tertentu dalam produksi. Aksi dijalankan
ketika posisi tersebut dicapai selama parsing.
\end{definisi}

\subsection{Contoh}

Untuk aturan
\[
E \to E_1 + E_2
\]
translation scheme dapat ditulis:
\[
E \to E_1 + \{ E.type = \text{arith}(E_1.type, E_2.type); \}\ E_2.
\]
Aksi dijalankan setelah \(E_1\) dan sebelum \(E_2\) diproses.

\subsection{Dari SDD ke Translation Scheme}

SDD S-attributed dapat diubah menjadi translation scheme dengan
menempatkan aksi di akhir produksi. SDD L-attributed memerlukan
penempatan aksi yang lebih hati-hati agar inherited attribute sudah
tersedia.

\begin{figure}[H]
\centering
\begin{tikzpicture}[node distance=6mm and 8mm]
  \node[phase, text width=4.4cm] (sdd) {SDD:\\aturan semantik};
  \node[phase, text width=4.4cm, right=of sdd] (scheme)
        {Translation scheme:\\aksi berurutan};
  \node[phase, text width=4.4cm, below=of scheme] (code)
        {Kode parser/visitor};
  \draw[arrow] (sdd) -- (scheme);
  \draw[arrow] (scheme) -- (code);
\end{tikzpicture}
\caption{Dari SDD ke translation scheme ke implementasi}
\label{fig:ch14-scheme}
\end{figure}

\subsection{Contoh pada Bison}

Semantic action Bison adalah bentuk translation scheme:
\begin{lstlisting}[language=Bison]
expr : expr T_PLUS expr
     {
         $$ = ast_new(AST_BINARY, $2.token);
         $$->left = $1;
         $$->right = $3;
     }
     ;
\end{lstlisting}
Aksi dijalankan saat aturan direduksi, setelah kedua operand
tersedia. Ini adalah pola S-attributed yang khas.

\subsection{Contoh pada Pernyataan}

Translation scheme juga berlaku untuk pernyataan. Assignment dan if
dapat ditulis sebagai:
\[
\begin{aligned}
S &\to \code{id} = E\ \{\, S.type =
\text{check\_assign}(\text{lookup}(\code{id}), E.type); \,\}\\
S &\to \code{if}\ (E)\ S_1\ \{\, \text{check\_bool}(E.type);\
S.type = S_1.type; \,\}
\end{aligned}
\]
Aksi assignment dijalankan setelah \(E\) selesai diproses, sehingga
tipe operand sudah tersedia. Aksi if memeriksa tipe kondisi sebelum
mewariskan tipe ke \(S\).

\subsection{Kapan Memilih}

SDD lebih mudah dibaca dan dianalisis. Translation scheme lebih
dekat ke implementasi dan mudah dipasang pada parser. Untuk
dokumentasi, SDD lebih baik; untuk implementasi, translation scheme
lebih praktis.

\begin{infobox}
Jangan mencampur logika parsing dan logika semantik yang kompleks
dalam satu aksi. Pindahkan perhitungan ke fungsi bantu agar grammar
tetap ringkas dan dapat diuji.
\end{infobox}
